\documentclass[11pt,a4paper]{article}

\usepackage[margin=1in]{geometry}
\usepackage{booktabs}
\usepackage{amsmath}
\usepackage{graphicx}
\usepackage{hyperref}
\usepackage{xcolor}
\usepackage{enumitem}
\usepackage{array}

\hypersetup{
    colorlinks=true,
    linkcolor=blue,
    urlcolor=blue,
    citecolor=blue
}

\title{\textbf{Grandmaster-RL: Adaptive Chess Puzzle Curriculum Learning}}
\author{DSAI Project -- IIIT Bangalore}
\date{}

\begin{document}

\maketitle

\begin{abstract}
Grandmaster-RL is a reinforcement learning based adaptive chess puzzle training system.
The system models chess puzzle recommendation as a sequential decision-making problem:
at each step, an agent selects a tactical theme and target puzzle difficulty for a
simulated learner. The environment retrieves real Lichess puzzles through per-theme
KD-tree indices and simulates solve outcomes using an Item Response Theory inspired
learner model. We compare Proximal Policy Optimization (PPO), Soft Actor-Critic (SAC),
and three heuristic baselines. PPO achieves the strongest overall learning efficiency,
with high solve rates and robust transfer to a held-out evaluation environment.
\end{abstract}

\section{Introduction}

Chess puzzle platforms commonly use rating-matching heuristics: a learner is given
puzzles near their current estimated rating. While this is simple and effective, it
does not explicitly optimize long-term skill development across tactical themes.
Grandmaster-RL instead treats adaptive puzzle selection as a Markov Decision Process
(MDP), where the agent learns to select curricula that maximize simulated learner
improvement over an episode.

The agent chooses both a tactical theme and a target difficulty. The environment then
retrieves the nearest real puzzle from a preprocessed Lichess puzzle dataset and
updates the learner's simulated skill based on whether the puzzle is solved. The final
goal is not merely high success rate, but efficient learning: rating gain achieved
through appropriately challenging and stable puzzle selection.

\section{Dataset and Preprocessing}

The project uses a reduced Lichess puzzle dataset. The preprocessing pipeline applies
the following steps:

\begin{enumerate}[leftmargin=*]
    \item Keep puzzles with rating deviation at most 150.
    \item Keep puzzles with ratings in the range $[400, 3000]$.
    \item Assign each puzzle to one primary tactical theme.
    \item Save one processed CSV per theme.
    \item Build a one-dimensional KD-tree per theme using puzzle rating as the search coordinate.
\end{enumerate}

One implementation issue was corrected during development. The internal theme name
\texttt{discovery} does not appear literally in the Lichess tags. Instead, Lichess
uses tags such as \texttt{discoveredAttack} and \texttt{discoveredCheck}. These tags
were mapped into the internal \texttt{discovery} category, producing a non-empty
discovery index.

\begin{table}[h]
\centering
\caption{Final per-theme puzzle counts after preprocessing.}
\begin{tabular}{lr}
\toprule
\textbf{Theme} & \textbf{Puzzle Count} \\
\midrule
Fork & 743,786 \\
Pin & 325,478 \\
Mate & 1,685,685 \\
Endgame & 1,576,401 \\
Skewer & 28,223 \\
Discovery & 162,344 \\
\bottomrule
\end{tabular}
\end{table}

\section{Environment Design}

The training environment is implemented as a Gymnasium environment. Each episode
contains 100 puzzle attempts. The observation vector contains:

\[
s_t =
\left[
\rho,\,
\phi_{\text{fork}},\,
\phi_{\text{pin}},\,
\phi_{\text{mate}},\,
\phi_{\text{endgame}},\,
\phi_{\text{skewer}},\,
\phi_{\text{discovery}},\,
\bar{p},\,
\bar{\delta}
\right],
\]

where $\rho$ is the overall learner rating, $\phi$ values are per-theme skill
estimates, $\bar{p}$ is rolling success rate, and $\bar{\delta}$ is rolling mean
difficulty. All observation values are normalized to $[0,1]$.

The action is a two-dimensional continuous vector:

\[
a_t = [k, d],
\]

where $k$ is a theme value rounded to the nearest valid theme index and $d$ is a
normalized target difficulty. The target difficulty is denormalized to the rating
range $[400, 3000]$, and the nearest real puzzle rating is retrieved from the
corresponding KD-tree.

\subsection{Learner Model}

The probability of solving a puzzle follows a logistic Item Response Theory style
model:

\[
P(\text{solve}) = \sigma\left(\frac{\phi_k - \delta}{\tau}\right),
\]

where $\phi_k$ is the learner's skill for the chosen theme, $\delta$ is the puzzle
rating, and $\tau$ controls the sharpness of the response curve.

After each attempt, the learner's skill is updated. Successful solves give a base
gain plus extra credit for solving puzzles above the learner's current skill, while
failures apply a fixed penalty:

\[
\text{gain} =
\begin{cases}
2\mu + \lambda \max(\delta - \phi, 0), & \text{if solved}, \\
-\mu, & \text{otherwise}.
\end{cases}
\]

\[
\phi_{\text{new}} = \text{clip}(\phi + \alpha \cdot \text{gain}, 400, 3000).
\]

\section{Algorithms}

\subsection{PPO}

PPO is the primary algorithm. It uses a custom hybrid actor-critic policy with:

\begin{itemize}[leftmargin=*]
    \item a categorical head for tactical theme selection,
    \item a Gaussian head for target difficulty,
    \item a shared MLP feature trunk,
    \item a value head for critic estimation.
\end{itemize}

PPO is well suited to the mixed discrete-continuous nature of the task because the
custom policy can explicitly model theme selection and difficulty selection separately.

\subsection{SAC}

SAC is included as a comparison model. Stable-Baselines3 SAC expects a fully continuous
policy, so the SAC trainer uses the standard \texttt{MlpPolicy}. This is compatible
with the environment because the action space is represented as a continuous
\texttt{Box}: the theme component is a float rounded inside the environment, and the
difficulty component is continuous. SAC is therefore useful as an aggressive
continuous-control comparison, though it is less semantically exact than PPO for theme
selection.

\section{Baselines}

Three non-RL baselines were used:

\begin{itemize}[leftmargin=*]
    \item \textbf{RandomAgent}: randomly selects theme and difficulty.
    \item \textbf{RatingMatchAgent}: selects difficulty near the learner's current rating.
    \item \textbf{FixedProgressionAgent}: starts with easier puzzles and increases difficulty on a fixed schedule.
\end{itemize}

These baselines are important because rating matching is already a strong practical
heuristic in puzzle recommendation systems.

\section{Evaluation Metrics}

The main metric is the Learning Efficiency Index (LEI):

\[
\text{LEI}
= \frac{\Delta \rho}{T}
\cdot \bar{p}
\cdot
\frac{1}{1 + \frac{\text{Var}(\delta)}{\sigma^2_{\text{ref}}}},
\]

where $\Delta \rho$ is total rating gain, $T$ is episode length, $\bar{p}$ is mean
success rate, and the final term rewards difficulty consistency.

We also report:

\begin{itemize}[leftmargin=*]
    \item mean rating change, $\Delta \rho$,
    \item mean success rate,
    \item robustness, computed as evaluation LEI divided by training LEI.
\end{itemize}

The evaluation environment adds unseen learner fatigue and response-temperature jitter.
This tests whether policies transfer beyond the exact training dynamics.

\section{Results}

\begin{table}[h]
\centering
\caption{Final comparison of RL agents and baselines.}
\resizebox{\textwidth}{!}{
\begin{tabular}{lrrrrrrr}
\toprule
\textbf{Agent} &
\textbf{LEI Train} &
\textbf{LEI Eval} &
\textbf{Robustness} &
\textbf{$\Delta \rho$ Train} &
\textbf{$\Delta \rho$ Eval} &
\textbf{Success Train} &
\textbf{Success Eval} \\
\midrule
PPO & 0.0024 & 0.0016 & 0.6674 & 6.38 & 4.63 & 0.754 & 0.697 \\
SAC & 0.0012 & 0.0005 & 0.4077 & 16.65 & 11.34 & 0.190 & 0.133 \\
Random & 0.0000 & 0.0000 & 0.6413 & 3.94 & 2.57 & 0.234 & 0.233 \\
RatingMatch & 0.0018 & 0.0012 & 0.6667 & 8.56 & 6.48 & 0.485 & 0.419 \\
FixedProgression & 0.0019 & 0.0013 & 0.7149 & 7.19 & 5.57 & 0.582 & 0.536 \\
\bottomrule
\end{tabular}
}
\end{table}

\section{Discussion}

PPO is the strongest overall policy according to the primary metric. It obtains the
highest evaluation LEI while maintaining a high success rate of 69.7\% in the held-out
evaluation environment. This indicates that PPO learns a curriculum that is both
challenging and consistently solvable.

SAC produces the largest raw rating gains, with an evaluation $\Delta \rho$ of 11.34.
However, its success rate is much lower: 19.0\% during training evaluation and 13.3\%
in the held-out environment. This suggests that SAC learns a more aggressive strategy,
selecting harder puzzles that produce large gains when solved but many failures overall.
Its lower robustness score also suggests weaker transfer to unseen learner dynamics.

The heuristic baselines are competitive. RatingMatch and FixedProgression both achieve
positive evaluation LEI, confirming that simple curriculum strategies are strong in this
domain. PPO still outperforms both on evaluation LEI, supporting the value of RL-based
curriculum optimization.

FixedProgression has the highest robustness score, but its evaluation LEI is below PPO.
This means it transfers consistently, but does not optimize learning efficiency as well
as the PPO policy.

\section{Implementation Fixes}

Several important issues were found and fixed during development:

\begin{itemize}[leftmargin=*]
    \item Corrected SciPy KD-tree query index extraction.
    \item Fixed empty \texttt{discovery.pkl} by mapping \texttt{discoveredAttack} and \texttt{discoveredCheck}.
    \item Added repository-root bootstrapping so training scripts run correctly as files.
    \item Disabled Stable-Baselines3 progress bars automatically when optional \texttt{rich} dependency is missing.
    \item Wrapped PPO callback evaluation environments with \texttt{VecNormalize}.
    \item Switched SAC from incompatible \texttt{HybridPolicy} to SB3's continuous \texttt{MlpPolicy}.
    \item Rebalanced the learner skill update so successful solves produce meaningful positive learning signal.
\end{itemize}

\section{Conclusion}

Grandmaster-RL successfully implements an adaptive chess puzzle curriculum system using
real puzzle metadata, simulated learner dynamics, and reinforcement learning. PPO is
the recommended final policy because it provides the best balance of learning
efficiency, success rate, rating improvement, and robustness. SAC is useful as an
aggressive high-gain comparison, while the heuristic baselines provide strong reference
points for practical curriculum design.

\section{Future Work}

Future improvements include implementing a true hybrid SAC policy with Gumbel-Softmax
theme selection, evaluating on more diverse learner profiles, adding richer theme
taxonomy, and validating the learned curricula with human users or historical learner
data.

\begin{thebibliography}{9}

\bibitem{glicko}
M. E. Glickman.
\textit{The Glicko System}.
Boston University, 1995.

\bibitem{lichess}
Lichess.
\textit{Lichess Open Puzzle Database}.
\url{https://database.lichess.org/#puzzles}

\bibitem{ppo}
J. Schulman, F. Wolski, P. Dhariwal, A. Radford, and O. Klimov.
\textit{Proximal Policy Optimization Algorithms}.
arXiv:1707.06347, 2017.

\bibitem{sac}
T. Haarnoja, A. Zhou, P. Abbeel, and S. Levine.
\textit{Soft Actor-Critic: Off-Policy Maximum Entropy Deep Reinforcement Learning with a Stochastic Actor}.
ICML, 2018.

\bibitem{gumbel}
E. Jang, S. Gu, and B. Poole.
\textit{Categorical Reparameterization with Gumbel-Softmax}.
ICLR, 2017.

\end{thebibliography}

\end{document}
